\section{Discussion}
\label{sec:discussion}
\subsection{NeuralNetRegressor}
\subsubsection{GridSearchCV}
One problem with the implementation of the GridSearchCV was the step wise
implementation, the first implmentation established 
the best method, the second implementation established the 
best activation function and hidden layers, and the third
established the penalty and the learning rate. If all
these values were tested in the same grid search, 
some other model parameter values may have been proven best,
as they vary according to the other model parameters. This
information was lost in the step wise implementation. What
was gained in this implementation is increased control,
data visualisation and shorter fit time per model, which made the implementation
more flexible. 
\subsubsection{Vanishing gradients}
The problem of vanishing gradients, when using sigmoid
in the hidden layers in deeper networks was not explored
sufficiently. There are two reasons for this, the first
is that the code updating the gradients, does not
include a method to store the gradients, such that one
could plot the evolving gradients, and compare them between the different layers.
One knows that it is the layers that are furthest from the 
output layer that are most affected by this, so an L2 norm
of the gradients in each layer could express, if the 
gradients were stored. \cite{VanishingGradient2024}
Also this project did not really explore networks with deep layers
as the deepest network included three hidden layers, this
could be explored in a seperate section by looking at the
gradients of the different layers, of a deeper network,
with sigmoid as activation function. The chosen network
in this project, based one the lowest one SE CV MSE, had one hidden layer, and was not affected
by this problem as strongly as a deeper network would. 
\subsubsection{Data cost}
The final test concluded that the neural network had the 
lowest test MSE, and by this criterion it was also the best
model. When looking at the fit time, the neural network does
not compare with the linear models, and the small difference
in test mse does not pay for the large difference in data cost.
This problem could be solved by code optimization, but was
not a part of this project. 

\subsection{NeuralNetclassifier}
For the neuralnetclassifier, and the GridSearchCV over
all model parameters, the data cost was the largest issue 
of this part of the project. The total data set was shortened into
1000 data points, 1000 data points out of a total 60 0000 training data.
This could lead one to think that the network would not 
have enough data to become sufficiently accurate, but an
accuracy of $89.22\%$ for the final model on the test data
shows that the model did learn, and classifies the unkown digits
with a rather high accuracy. 
Adding a simple logistic regression model, and 
fitting the model on the same data, would have further shown
how well the network actually performed, and this is a 
problem with this classification problem, it lacks models
for comparison, as we had with the regression problem.
 

